\documentclass[11pt]{article}
\usepackage{nopageno}
\setcounter{secnumdepth}{1}
\usepackage[margin=1in]{geometry}
\usepackage{titlesec}
\titleformat{\section}[runin]{\normalfont\bfseries}{\thesection}{1em}{}[]
\titlespacing*{\section}{0pt}{0.1\baselineskip}{1ex}
\titleformat{\subsection}[runin]{\normalfont\bfseries}{\thesubsection}{1em}{}[:]
\titlespacing*{\subsection}{0pt}{0.1\baselineskip}{1ex}
\usepackage{enumitem}
\begin{document}
\pagestyle{empty}

\noindent \textbf{Data Management and Sharing Plan}
\vspace{0.3cm}

\section{Roles and Responsibilities}
The Principal Investigator (PI), Dr.\ Rebecca Napolitano, will have
primary responsibility for data management throughout the project lifecycle and will
ensure data is collected, stored, preserved, and shared in accordance with this plan
and NSF policies.
Co-PI Wauthier is responsible for SAR data products, including ICEYE-derived
building-level inundation measurements and Sentinel-1 flood-extent imagery.
Co-PI Busse is responsible for hydrological model inputs, outputs, and associated
metadata.
All three PIs are at Penn State.

\vspace{0.3cm}
\section{Expected Data and Formats}
This project will generate the following data types:

\begin{itemize}
\item[] \textbf{ICEYE commercial SAR data (restricted):}
  Building-level flood depth and inundation duration measurements for 271 downtown
  Montpelier buildings across the 2023 and 2024 flood events, acquired under an NDA
  data access agreement with ICEYE.
  Raw ICEYE imagery and building-level measurements cannot be redistributed under
  the data access agreement; derived products (per-building duration statistics,
  validation metrics, anonymized scatter plots) will be shared.
  Formats: CSV summary tables (.csv), PDF figures (.pdf).
\item[] \textbf{Sentinel-1 SAR flood-extent products:}
  Amplitude imagery and flood-extent classifications for Montpelier, VT across both
  flood events, used for broader spatial context and cross-checking of ICEYE
  building-level classifications.
  Formats: GeoTIFF (.tif), CSV summary tables (.csv).
\item[] \textbf{NOAA precipitation and forecast inputs:}
  NOAA precipitation data and forecast products used for pre-storm staging and
  mass balance model inputs.
  Formats: CSV (.csv), NetCDF (.nc).
\item[] \textbf{Hydrological mass balance model outputs:}
  Drainage timescale estimates, sub-basin parameterization, watershed routing
  outputs, and uncertainty envelopes for the Winooski watershed.
  Formats: NetCDF (.nc), CSV (.csv).
\item[] \textbf{Bayesian posterior distributions:}
  Posterior duration distributions per building, credible intervals, and error
  model parameters.
  Formats: NetCDF (.nc), CSV summary tables (.csv).
\item[] \textbf{Fragility surface parameters:}
  Duration-depth fragility surfaces for pre-code load-bearing masonry archetypes,
  calibrated on 2023 observations and validated against 2024.
  Formats: CSV (.csv), Python serialization (.pkl).
\item[] \textbf{Field damage observations:}
  Building-level damage state classifications (no damage, cosmetic, structural)
  from field assessment of 271 downtown Montpelier buildings.
  These data are building-identifiable; anonymized versions with parcel-level
  identifiers removed will be shared publicly, and the full geo-referenced dataset
  will be available to qualified researchers under a data use agreement.
  Formats: CSV (.csv), GeoJSON (.geojson).
\item[] \textbf{Community co-development workshop materials:}
  Workshop documentation, structured elicitation results, participant feedback,
  and facilitation guides.
  Formats: PDF (.pdf), anonymized transcripts (.txt, .docx).
\item[] \textbf{Decision framework outputs:}
  Retrofit priority rankings, what-if simulation outputs, and community-scale
  drainage versus individual building retrofit comparisons.
  Formats: CSV (.csv), PDF reports for community partners.
\item[] \textbf{Dissemination materials:}
  Publications, presentations, webinars (.pdf, .pptx, .mp4).
\end{itemize}

Metadata will include instrument specifications (SAR sensor, orbit, acquisition dates),
data processing parameters, model version information, and field data collection
protocols.

\vspace{0.3cm}
\section{Access, Sharing, and Reuse Policies}
SAR-derived flood-extent products, mass balance outputs, posterior distributions,
and fragility surface parameters will be shared via DesignSafe-CI under open-access
licenses upon publication or within 12 months of data collection completion.
ICEYE commercial SAR data is acquired under a data access agreement that restricts
redistribution of raw imagery; derived per-building statistics and validation
metrics will be shared on DesignSafe-CI.
Bayesian fusion, mass balance parameterization, and decision framework code will
be released on GitHub under an MIT open-source license.
Field damage observations will be shared on the Open Science Framework: anonymized
data publicly, geo-referenced data under a data use agreement.
Community workshop materials will be anonymized before sharing; access to
identifiable materials will be restricted to project personnel.
Publications and presentations will be made openly available via institutional
repositories.

The hazard-related data and metadata will be archived on the Data Depot at
DesignSafe-CI (the standard repository for the natural hazards community), following
established standards listed on the DesignSafe-CI website and subject to
DesignSafe-CI Data Sharing and Archiving Policies.

\section{Reproducibility}
All data processing and analysis scripts will be version-controlled using Git and
stored on GitHub with appropriate documentation.
Computational environments will be documented using Docker containers or Conda
environment files, with configuration files shared alongside data.
All statistical analyses and visualizations will include code to reproduce them from
processed data.
Open-source software tools will be prioritized to remove licensing barriers.

\section{Data Storage, Preservation and Backup}
All data and code will be archived on DesignSafe-CI.
PI Napolitano is an affiliate of Penn State's Institute for Computation and Data
Science (ICDS) and, through its Advanced Cyber Infrastructure (ICDS-ACI), has access
to active and near-line storage that is backed up daily.
Co-PIs Wauthier and Busse will use ICDS-ACI storage for active research data,
ensuring all project data resides within Penn State's backed-up infrastructure.

\vspace{0.3cm}
\section{Ethics and Privacy}
The project will comply with all relevant regulations, including Institutional Review
Board (IRB) requirements for human subjects research.
Community co-development workshop participants will be fully informed of data
collection and sharing plans.
Workshop notes will be anonymized before sharing; personally identifiable information
will be protected using encryption and access controls.
Field damage observations are building-identifiable; publicly shared versions will
have parcel-level identifiers removed, and geo-referenced data will be available only
under a data use agreement to protect building owner privacy.
This plan will be reviewed annually and revised as needed in consultation with the
Penn State IRB.

\end{document}
